%% 01_introduccion.tex — Contexto, desafíos y contribuciones.

\section{Introducción}
\label{sec:intro}

El análisis dinámico de seguridad de aplicaciones (DAST) evalúa un sistema en
ejecución enviando peticiones manipuladas y observando la respuesta, sin
acceso al código fuente. Frente al análisis estático, su principal ventaja es
la ausencia de suposiciones sobre el lenguaje o el marco de trabajo del
objetivo; su principal inconveniente es el coste. Diversos estudios
comparativos han mostrado que los escáneres de caja negra presentan cobertura
desigual y tasas de falsos positivos
elevadas~\cite{doupe2010johnny,bau2010state,parvez2015analysis}, lo que erosiona
la confianza del equipo de desarrollo y, en la práctica, conduce a ignorar los
informes.

Identificamos tres desafíos concretos que este trabajo aborda de forma
explícita:

\begin{enumerate}
  \item \textbf{Latencia acumulada.} Un escaneo realista implica del orden de
  $10^4$--$10^5$ peticiones HTTP. Ejecutadas secuencialmente, la latencia de
  red domina el tiempo total y hace inviable integrar el DAST en cada
  \emph{commit}.
  \item \textbf{Ruido (falsos positivos).} La detección basada únicamente en
  la reflexión de una cadena o en un patrón de error de base de datos produce
  hallazgos que no son explotables. Distinguir una reflexión inerte de una
  ejecución real requiere una segunda observación dirigida.
  \item \textbf{Explosión combinatoria.} El producto (puntos de inyección)
  $\times$ (\emph{payloads} por categoría) $\times$ (transformaciones de
  evasión) crece hasta hacer imposible probarlo todo bajo un presupuesto de
  tiempo fijo; el orden en que se agota ese presupuesto determina qué se
  detecta.
\end{enumerate}

\noindent\textbf{Contribuciones.} Este artículo realiza las siguientes
aportaciones:

\begin{itemize}
  \item Un motor DAST asíncrono con un único punto de estrangulamiento HTTP,
  control de concurrencia con contrapresión, limitación de tasa por
  \emph{token bucket} y salvaguardas frente a SSRF en el seguimiento de
  redirecciones (\cref{sec:arquitectura}).
  \item Un \emph{pipeline} de planificación de \emph{payloads} con etapas
  desacoplables --- entrelazado por contexto, priorización y ordenación --- y
  una etapa de confirmación en tiempo de ejecución que calibra la confianza
  del hallazgo (\cref{sec:arquitectura,sec:metodologia}).
  \item Un \code{TelemetryWorker} no bloqueante que registra una fila JSONL
  por \emph{sonda} de \emph{payload}, habilitando la medición del coste de
  detección instancia a instancia (\cref{sec:telemetria}).
  \item Una metodología de evaluación reproducible: banco de pruebas
  contenedorizado, oráculo de calificación independiente, diseño experimental
  RCBD y manifiesto de reproducibilidad con \emph{hash} del corpus
  (\cref{sec:metodologia}).
  \item Un estudio de ablación que aísla la contribución de cada etapa del
  \emph{pipeline} a la curva de detección frente a coste y a la tasa de
  falsos positivos, y que revela que la palanca de eficiencia no es la etapa
  concreta sino el \emph{ajuste del presupuesto} de peticiones: un punto óptimo
  de $F_1$ bajo escasez, degradación de la precisión por saturación y un techo
  de detección atribuible a la cobertura del corpus (\cref{sec:resultados}).
\end{itemize}

El resto del artículo se organiza como sigue. La \cref{sec:arquitectura}
describe el diseño del motor. La \cref{sec:metodologia} detalla el banco de
pruebas y el diseño experimental. La \cref{sec:resultados} presenta la
evaluación. La \cref{sec:amenazas} discute las amenazas a la validez y la
\cref{sec:conclusion} concluye.
